\documentclass{fisatprojectfinal}
\renewcommand{\baselinestretch}{1.50} 
\usepackage{setspace}
\usepackage[utf8]{inputenc}
\usepackage{xcolor}
\usepackage{enumitem}
\usepackage{tabularx}
\usepackage{longtable}
\usepackage{booktabs}
\usepackage{float}
\usepackage{appendix}

% Header and Footer setup
\usepackage{fancyhdr}
\pagestyle{fancy}
\fancyhf{}
\fancyhead[L]{\textbf{APISEC}}
\fancyfoot[C]{\textbf{Computer Science \& Engineering --- FISAT}}
\fancyfoot[R]{\thepage}
\renewcommand{\headrulewidth}{0.6pt}
\renewcommand{\footrulewidth}{0.6pt}

\fancypagestyle{plain}{
  \fancyhf{}
  \fancyhead[L]{\textbf{APISEC}}
  \fancyfoot[C]{\textbf{Computer Science \& Engineering --- FISAT}}
  \fancyfoot[R]{\thepage}
  \renewcommand{\headrulewidth}{0.6pt}
  \renewcommand{\footrulewidth}{0.6pt}
}

\title{APISEC - AI-Enhanced API Schema Monitor and Security Analyzer \\ \large Final Project Report}
\team{Nikhil (FITCS20001) \\ Teammate1 Name (FITCS20002) \\ Teammate2 Name (FITCS20003) \\ Teammate3 Name (FITCS20004)}
\author{Nikhil}
\regno{FITCS20001}
\bibliographystyle{plain}

\begin{document}
\pagestyle{empty}
\maketitle
\makecert
\newpage

\pagenumbering{roman}
\setcounter{page}{1}
\pagestyle{plain}

\newgeometry{top=4cm,bottom=0.1cm}
\renewcommand\abstractname{ABSTRACT}
\begin{abstract}
\vspace{2cm}
\begin{onehalfspace}
APISEC is an innovative, AI-driven platform designed to address the critical challenges of API governance, schema drift detection, and security analysis in modern microservice architectures. As contemporary software ecosystems increasingly rely on interconnected APIs, the lack of automated, semantically-aware monitoring tools often leads to undetected breaking changes, integration failures, and security vulnerabilities categorized under OWASP's Improper Assets Management. APISEC introduces a comprehensive framework for autonomous schema discovery, immutable versioning, and intelligent differential analysis.

The system is architected using a modern full-stack approach, utilizing React 19 with Vite for a high-performance frontend, FastAPI (Python) for an asynchronous backend, and PostgreSQL for robust multi-tenant data persistence. A key innovation of APISEC is its integration with the Google Gemini \texttt{gemini-2.5-flash} model, which translates complex structural deltas into human-readable impact reports and actionable remediation guidance. Furthermore, the platform incorporates a concurrent Runtime Validator for live contract conformance checking and a Differential Probing Engine that utilizes behavioral fingerprinting to uncover undocumented API parameters.

Our implementation demonstrates that the integration of generative AI into API governance workflows significantly reduces the cognitive overhead for engineering teams. The system successfully isolates semantic schema changes from cosmetic reformatting, providing a reliable audit trail and a proactive security net. This report details the design, algorithms, and implementation of APISEC, showing its effectiveness in maintaining API contract integrity and enhancing developer productivity across the software development lifecycle.
\end{onehalfspace}
\end{abstract}
\newpage

\renewcommand\abstractname{Contribution by Author}
\begin{abstract}
\vspace{5cm}
\textbf{Nikhil} engineered the native AI integration and differential logic, constructing the core schema mapping processes querying Google Gemini \texttt{gemini-2.5-flash}. This included building the \texttt{ai\_analyzer.py} module with semantic prompt composition via \texttt{\_build\_prompt} and \texttt{analyze\_single\_change}, implementing the \texttt{\_summarize\_schema} truncation strategy for large nested definitions, and integrating the Differential Engine output into structured Gemini prompts producing human-readable impact reports and remediation guidance for breaking changes.\\[6pt]
\textbf{Teammate1 Name} designed and implemented the complete frontend ecosystem using React 19 and Vite. This included building responsive interfaces for the API Dashboard, the Schema Monitor with advanced search capabilities, on-demand Analyze Diff AI buttons, and JWT-based session management covering Login and Register views. Styling was achieved with TailwindCSS, PostCSS, and Lucide React icon components.\\[6pt]
\textbf{Teammate2 Name} architected the PostgreSQL database layer and FastAPI backend routing, implementing the \texttt{registry\_db.py} raw SQL schema managing \texttt{users}, \texttt{apis}, and \texttt{schema\_snapshots} tables with cascading foreign key constraints. She enforced multi-tenant data isolation through JWT bearer token validation on all core endpoints, designed the \texttt{/auth} registration and login flows, and integrated ReportLab for base64-encoded schema PDF generation on each new snapshot version.\\[6pt]
\textbf{Teammate3 Name} developed the runtime validation and probing infrastructure, implementing \texttt{runtime\_validator.py} with the \texttt{\_test\_endpoint} constraint generator and the \texttt{\_validate\_response\_schema} auditor. She also engineered the \texttt{probes/differential\_engine.py} module encompassing \texttt{ResponseFingerprint}, \texttt{StringProbe}, \texttt{NumericProbe}, and \texttt{BooleanProbe} classes, and configured the Docker Compose deployment stack defining the \texttt{apisec} backend container and \texttt{nginx} reverse proxy service.
\vspace{1cm}
\begin{flushright}
Nikhil\\
Teammate1 Name\\
Teammate2 Name\\
Teammate3 Name
\end{flushright}
\end{abstract}
\newpage

\renewcommand\abstractname{ACKNOWLEDGMENT}
\begin{abstract}
\vspace{5cm}
We express our heartfelt gratitude to our Principal, for providing an inspiring academic environment and unwavering support that greatly contributed to the success of this project. We are also deeply thankful to our Head of the Department, Department of Computer Science Engineering, Federal Institute of Science and Technology (FISAT), for their constant encouragement and leadership. Additionally, we extend our sincere appreciation to our project coordinators for their valuable guidance, timely feedback, and dedicated mentorship throughout the project.

We would like to express our sincere gratitude to our project guide for invaluable guidance, constant encouragement, and constructive feedback throughout the duration of this project. His expertise and insights were instrumental in shaping our work.

We are also thankful to the Department of Computer Science Engineering, Federal Institute of Science and Technology (FISAT), for providing us with the necessary resources and infrastructure to carry out this project.

Finally, we would like to thank our families and friends for their unwavering support and encouragement during this academic endeavour.
\end{abstract}
\newpage

\restoregeometry
\tableofcontents
\newpage
\listoffigures
\newpage
\listoftables
\newpage
\chapter*{Glossary and Abbreviations}
\addcontentsline{toc}{chapter}{Glossary and Abbreviations}
\begin{description}[style=multiline,leftmargin=3cm,font=\bfseries]
    \item[API] Application Programming Interface
    \item[REST] REpresentational State Transfer
    \item[JSON] JavaScript Object Notation
    \item[YAML] YAML Ain't Markup Language
    \item[JWT] JSON Web Token
    \item[OWASP] Open Web Application Security Project
    \item[SRS] Software Requirement Specification
    \item[HTML] HyperText Markup Language
    \item[SPA] Single Page Application
    \item[ASGI] Asynchronous Server Gateway Interface
    \item[CRUD] Create, Read, Update, Delete
    \item[LLM] Large Language Model
    \item[NLP] Natural Language Processing
    \item[SDK] Software Development Kit
    \item[SSL] Secure Sockets Layer
    \item[HTTP] HyperText Transfer Protocol
\end{description}
\newpage

% STARTING INTRO - HEADER/FOOTERS MUST BE APPLIED
\pagenumbering{arabic}
\setcounter{page}{1}
\pagestyle{fancy}

\chapter{Introduction}
\section{Overview}
In the era of cloud-native computing and microservice architectures, Application Programming Interfaces (APIs) have become the fundamental building blocks of modern software. They facilitate seamless communication between disparate systems, enabling the complex workflows that power today's digital economy. However, this increased reliance on APIs has introduced a new class of challenges related to contract maintenance, schema evolution, and security governance.

APISEC (API Security and Schema Monitor) is presented as a next-generation solution to these challenges. Traditional methods of tracking API changes often rely on manual documentation updates or textual diffs generated by version control systems like Git. These methods are fundamentally flawed when applied to structured API specifications such as OpenAPI or Swagger documents. A textual diff may flag a massive number of changes due to simple key reordering or formatting updates, while simultaneously missing a critical, breaking change buried deep within a nested JSON structure.

APISEC addresses this by implementing a semantically-aware Differential Engine that understands the hierarchy and constraints of API contracts. By autonomously crawling for schemas, normalizing them into a canonical form, and storing versioned snapshots, APISEC provides a clear and immutable history of an API's evolution. The addition of Google Gemini \texttt{gemini-2.5-flash} takes this analysis a step further, providing developers with natural language explanations of changes, impact assessments for downstream consumers, and precise fix suggestions. This technical report explores the architecture, implementation, and evaluation of APISEC as a pivotal tool for ensuring API reliability and security.

\section{Problem Statement}
The rapid pace of modern software development often results in "API Drift," where the live implementation of an API diverges from its documented schema. This drift is frequently silent and remains undetected until it causes a runtime failure in a production environment. Such failures are not only costly to debug but also pose significant security risks. The OWASP API Security Top 10 highlights "Improper Assets Management" as a core vulnerability, often stemming from the use of outdated or undocumented API versions.

Current development workflows lack a unified tool that can:
\begin{enumerate}
    \item \textbf{Autonomously Discover} schemas across diverse environments without manual registration of every endpoint.
    \item \textbf{Differentiate Semantically} between structural changes and purely cosmetic modifications in JSON/YAML specifications.
    \item \textbf{Communicate Impact} in a way that is understandable to human developers, rather than just outputting raw structural deltas.
    \item \textbf{Validate Conformance} of live runtime responses against the published schema in a concurrent and scalable manner.
    \item \textbf{Expose Hidden Attack Surfaces} such as undocumented parameters that may be vulnerable to injection or data exposure.
\end{enumerate}
The absence of such a comprehensive governance platform leads to engineering bottlenecks, fragile integrations, and an expanded threat landscape for API-driven organizations.

\section{Objectives}
The primary objective of this project is to develop and implement APISEC, a comprehensive API governance and security analysis platform. The specific technical goals are:
\begin{enumerate}
    \item \textbf{Cloud-Native Dashboard}: To design a responsive React 19 frontend that provides a centralized view for API registration, scan history, and AI-powered diff analysis.
    \item \textbf{Canonical Normalization Engine}: To engineer a system that resolves \texttt{\$ref} pointers and sorts schema keys deterministically to eliminate false positives in differential analysis.
    \item \textbf{AI-Enriched Differential Engine}: To integrate the Gemini 2.5 Flash model to provide semantic context, impact analysis, and remediation guidance for detected schema variations.
    \item \textbf{Asynchronous Runtime Validator}: To build a high-performance auditor using \texttt{aiohttp} that probes endpoints concurrently to ensure live compliance with the defined schema contract.
    \item \textbf{Behavioral Fingerprinting Engine}: To develop a probing system capable of discovering undocumented request parameters through recursive injection and error-pattern analysis.
    \item \textbf{Multi-Tenant Persistence Layer}: To implement a secure PostgreSQL backend that manages versioned schema snapshots and user-specific registries with JWT authentication.
\end{enumerate}

\section{Scope of the Project}
The scope of APISEC encompasses the entire lifecycle of API schema governance, from discovery to runtime validation. The project focuses specifically on RESTful APIs and supports both OpenAPI (v2 and v3) and Swagger specifications. The technical implementation includes:
\begin{itemize}
    \item A FastAPI-based backend capable of handling high-concurrency tasks such as schema crawling and the concurrent probing of hundreds of endpoints.
    \item A React 19 frontend utilizing modern hooks and state management for an interactive user experience.
    \item Integration with the Google Gemini API for advanced NLP-based schema analysis.
    \item Support for PostgreSQL as the primary data store for snapshots and user data.
    \item Deployment configurations using Docker Compose and Nginx for standardized, containerized environments.
\end{itemize}
The project excludes active exploitation tools or penetration testing frameworks, focusing instead on contract integrity, structural security, and governance auditing.

\section{Social and Technical Relevance}
APISEC delivers significant value to the broader software engineering community and the digital economy:
\begin{itemize}
    \item \textbf{Enhanced Productivity}: Automating the tedious task of API contract auditing allows developers to focus on feature delivery rather than debugging integration mismatches.
    \item \textbf{Improved Security Posture}: By identifying undocumented endpoints and flagging missing authentication headers, APISEC directly reduces the risk of data breaches.
    \item \textbf{Operational Stability}: Preventing breaking changes from reaching production ensures the reliability of critical social infrastructure services that depend on APIs.
    \item \textbf{Governance and Compliance}: The immutable snapshot history and PDF exports provide a robust audit trail for organizations following strict regulatory frameworks (e.g., GDPR, HIPAA).
\end{itemize}

\section{Organization of the Report}
This report is organized as follows:
\begin{itemize}
    \item \textbf{Chapter 1}: Provides an overview of the project, problem statement, and objectives.
    \item \textbf{Chapter 2}: Reviews existing literature and positions APISEC against current industry tools.
    \item \textbf{Chapter 3}: Details the system design, including architecture, database schema, and core algorithms.
    \item \textbf{Chapter 4}: Explains the implementation details, technology stack, and deployment strategy.
    \item \textbf{Chapter 5}: Describes the testing methodologies and presents the results of the QA process.
    \item \textbf{Chapter 6}: Summarizes the project outcomes and provides a comparative analysis.
    \item \textbf{Chapter 7}: Concludes the report and outlines future research directions.
\end{itemize}


\chapter{Literature Review}

\section{Related Work}
The landscape of API monitoring, testing, and governance has evolved significantly with the rise of microservices. Several established tools provide various aspects of the functionality offered by APISEC.

\textbf{Postman} \cite{postman} is perhaps the most ubiquitous tool in the API developer's toolkit. It provides a robust environment for building, testing, and documenting APIs. Postman excels at request-level testing and workflow automation. However, its schema management is largely a premium feature and often requires manual intervention. It lacks an autonomous "crawler" that can discover schemas in the wild and does not currently integrate generative AI for semantic impact analysis of schema changes.

\textbf{Swagger and OpenAPI Tools} \cite{swagger} form the standard foundation for modern API documentation. Tools like Swagger UI and Swagger Editor allow developers to visualize and interact with their schemas. While excellent for documentation, these tools are static in nature. They do not provide versioned snapshot storage or an automated engine to compare two versions of a live schema to find breaking changes.

\textbf{Akita} \cite{akita} takes a unique approach by observing network traffic to build a dynamic model of an API. This "API Observability" allows Akita to detect schema drift in real-time. However, Akita requires integration at the infrastructure level (e.g., eBPF or proxy and can be complex to set up. APISEC, by contrast, uses a "shift-left" approach of crawling specifications directly, which is less intrusive and easier to deploy in diverse environments.

\textbf{42Crunch} \cite{fortytwo} is a high-end API security platform that focuses on the entire API lifecycle. It provides excellent schema auditing and runtime protection. While highly capable, it is targeted at enterprise-level security teams and can be cost-prohibitive for smaller development groups. APISEC provides a more focused, developer-centric alternative with a specific emphasis on AI-driven change interpretation.

\textbf{Spectral} \cite{spectral} is an open-source JSON/YAML linter that can be used to enforce organizational standards on OpenAPI documents. While it helps in maintaining schema quality, it does not address the problem of evolution over time or the discovery of undocumented endpoints.

\section{Comparison of Related Works}

\begin{table}[H]
\centering
\caption{Detailed Comparison of API Governance Tools}
\label{tab:comparison_detailed}
\begin{tabularx}{\textwidth}{|l|X|X|X|X|}
\hline
\textbf{Feature} & \textbf{Postman} & \textbf{Swagger Hub} & \textbf{Akita} & \textbf{APISEC} \\
\hline
\textbf{Discovery} & Manual & Static Registry & Traffic Observation & Autonomous Crawl \\
\hline
\textbf{Versioning} & Manual & Manual/Git & Automatic & Automated Snapshots \\
\hline
\textbf{Diff Type} & Textual & Semantic & Security-Focused & Structural + Semantic \\
\hline
\textbf{AI Analysis} & Limited & No & No & Gemini 2.5 Flash \\
\hline
\textbf{Validation} & Manual Tests & Static Lint & Runtime Observation & Concurrent Probing \\
\hline
\textbf{Accessibility} & Broad & Enterprise & Infrastructure-heavy & Developer-centric \\
\hline
\end{tabularx}
\end{table}

\section{Recent Trends in API Security}
The API security landscape is currently undergoing a paradigm shift, driven by several key trends that APISEC is designed to address.

\textbf{Shift-Left Security}: There is a growing emphasis on integrating security testing earlier in the software development lifecycle (SDLC). By analyzing schemas and detecting drift during development and staging, APISEC enables teams to catch vulnerabilities before they reach production, reducing remediation costs and risks.

\textbf{AI-Enriched Governance}: As APIs become more complex, manual auditing is no longer feasible. The use of LLMs for semantic analysis, as implemented in APISEC, represents a major trend toward "Self-Healing" or "Self-Documenting" API ecosystems where AI assists in explaining and fixing breaking changes.

\textbf{Shadow API Discovery}: Undocumented or "Shadow" APIs are a leading cause of data breaches. Modern tools are increasingly focusing on autonomous discovery mechanisms—like APISEC's multi-path crawler and behavioral fingerprinting—to uncover and secure hidden attack surfaces.

\section{Positioning of APISEC}
APISEC occupies a unique niche in the API toolchain. Unlike observational tools (Akita) that require complex infrastructure integration, APISEC is highly portable and can be deployed as a sidecar or a standalone auditor. It differentiates itself from standard documentation tools (Swagger) by being dynamic and version-aware. The most significant innovation is the application of Large Language Models (LLMs) to the problem of schema drift. While other tools may show a developer that a field name has changed from \texttt{userId} to \texttt{user\_id}, APISEC's Gemini integration can explain that this change will break the mobile app's core login flow and requires a specific update in the authentication middleware. This "semantic gap" closure is what positions APISEC as a transformative tool for modern engineering teams.


\chapter{Design Methodologies}

\section{Software Requirement Specifications}
\subsection{Functional Requirements}
APISEC shall provide the following functionality:

\begin{enumerate}
    \item \textbf{User Authentication and Multi-Tenancy}
    \begin{itemize}
        \item Secure user registration storing bcrypt-hashed credentials in PostgreSQL.
        \item Login endpoint issuing signed JWT access tokens.
        \item All API registry and schema snapshot endpoints protected by JWT bearer token validation enforcing per-user data isolation.
    \end{itemize}

    \item \textbf{API Registry Management}
    \begin{itemize}
        \item Create, list, and delete API targets storing name, base URL, description, and owner user ID.
        \item Cascade deletion of all associated schema snapshots when an API entry is removed.
    \end{itemize}

    \item \textbf{Schema Discovery and Crawling}
    \begin{itemize}
        \item Autonomous crawling of registered API base URLs by appending common OpenAPI/Swagger spec paths (\texttt{/openapi.json}, \texttt{/swagger.json}, \texttt{/docs}, etc.) in priority order.
        \item HTML fallback parsing via \texttt{discover\_schema\_from\_html} extracting \texttt{<link>} and \texttt{<script>} references to embedded schemas.
        \item Ad-hoc schema discovery via \texttt{POST /discover-schema} for unauthenticated preview.
    \end{itemize}

    \item \textbf{Schema Versioning and Snapshot Storage}
    \begin{itemize}
        \item Incremental version numbering per API; new snapshot created only when a structural change is detected against the latest stored version.
        \item Schema stored as serialized JSON and as a base64-encoded PDF generated by ReportLab.
        \item Endpoints to retrieve full snapshot history, fetch the latest version, or retrieve a specific version by number.
    \end{itemize}

    \item \textbf{Schema Normalization and Differential Analysis}
    \begin{itemize}
        \item Normalization Engine stripping non-contractual metadata, recursively resolving \texttt{\$ref} imports, and deterministically sorting keys before comparison.
        \item Differential Engine traversing full JSON trees, classifying property changes as added, removed, or modified with severity ratings.
        \item Structured diff output selectable via \texttt{?structured=True} query parameter.
    \end{itemize}

    \item \textbf{AI-Powered Change Analysis}
    \begin{itemize}
        \item On-demand invocation of Google Gemini \texttt{gemini-2.5-flash} for a specific schema diff, producing: (1) human-readable change description, (2) consumer and SDK impact analysis, (3) actionable remediation steps for breaking changes.
    \end{itemize}

    \item \textbf{Runtime Validation}
    \begin{itemize}
        \item Concurrent asynchronous probing of all schema-defined endpoints using \texttt{aiohttp} with semaphore-bounded concurrency.
        \item Deterministic test parameter generation seeded by method, path, and parameter name for repeatable results.
        \item Response auditing comparing actual HTTP status codes and JSON payload shapes against schema-defined expectations.
    \end{itemize}

    \item \textbf{Differential Probing Engine}
    \begin{itemize}
        \item Behavioral fingerprinting via \texttt{StringProbe}, \texttt{NumericProbe}, and \texttt{BooleanProbe} injected recursively into endpoint request bodies.
        \item Extraction of hidden parameter names from error messages and behavioral variance hashing.
    \end{itemize}
\end{enumerate}

\subsection{Non-Functional Requirements}
\begin{itemize}
    \item \textbf{Performance}: Schema diff operations shall complete within sub-second latency for typical OpenAPI documents. Runtime validation shall leverage \texttt{aiohttp} concurrency to avoid blocking the FastAPI event loop.
    \item \textbf{Security}: User credentials shall be stored using secure hashing. All data access shall be gated behind JWT bearer token validation enforcing strict per-user data isolation.
    \item \textbf{Scalability}: The PostgreSQL schema shall support horizontal scaling of API registries and snapshot histories without structural changes. Docker Compose orchestration allows independent scaling of the backend and proxy layers.
    \item \textbf{Reliability}: The Schema Crawler shall implement a multi-path fallback strategy ensuring maximum schema discovery success even for non-standard API deployments.
    \item \textbf{Maintainability}: Core analytical modules (Schema Monitor, AI Analyzer, Runtime Validator, Differential Probing Engine) shall be implemented as independent Python modules with clearly defined interfaces.
\end{itemize}

\section{Software Design Document}
\subsection{System Features}
The system is organised into five core functional modules:
\begin{itemize}
    \item \textbf{Authentication Module}: JWT-secured user registration and login enforcing multi-tenant data isolation.
    \item \textbf{API Registry Module}: CRUD management of monitored API targets with per-user ownership.
    \item \textbf{Schema Monitor Module}: Autonomous crawling, normalization, versioning, and differential analysis of OpenAPI schemas.
    \item \textbf{AI Analyzer Module}: Google Gemini-powered semantic interpretation of schema diffs.
    \item \textbf{Runtime and Probing Module}: Concurrent endpoint validation and behavioral fingerprinting for hidden parameter discovery.
\end{itemize}

\subsection{System Architecture Design}

\begin{figure}[h!]
    \begin{center}
        \includegraphics[width=0.8\textwidth]{placeholder}
        \caption{APISEC System Architecture Diagram}
    \end{center}
\end{figure}

The APISEC architecture is structured into four key layers. The \textbf{Client Layer} consists of a React 19 single-page application built with Vite, styled using TailwindCSS, and communicating with the backend via authenticated REST requests carrying JWT bearer tokens. The \textbf{Application Layer} is a Python FastAPI application (\texttt{main.py}) running on Uvicorn, acting as the central hub connecting the frontend to the core analytical modules. All sensitive endpoints are protected by JWT middleware. The \textbf{Core Modules Layer} contains the independently deployable Python modules: \texttt{schema\_monitor.py}, \texttt{ai\_analyzer.py}, \texttt{runtime\_validator.py}, and \texttt{probes/differential\_engine.py}. The \textbf{Data Layer} comprises a PostgreSQL relational database accessed via raw SQL through \texttt{registry\_db.py}, storing user accounts, API registries, and versioned schema snapshots. The entire stack is containerised using Docker Compose, with an \texttt{nginx} container handling reverse proxy duties in front of the FastAPI backend.

\subsection{Database Design}
The system uses PostgreSQL with a simple relational schema managed via raw SQL queries in \texttt{registry\_db.py}. Three primary tables form the data backbone:

The \textbf{users} table manages authenticated identities for multi-tenant data isolation. It contains \texttt{id} (SERIAL PRIMARY KEY), \texttt{username} (TEXT UNIQUE), \texttt{password\_hash} (TEXT), and \texttt{created\_at} (TIMESTAMP).

The \textbf{apis} table tracks registered API monitoring targets, isolated per user. Key columns are \texttt{id} (SERIAL PRIMARY KEY), \texttt{user\_id} (INTEGER, foreign key to \texttt{users.id} with ON DELETE CASCADE), \texttt{name} (TEXT), \texttt{base\_url} (TEXT), \texttt{description} (TEXT), and \texttt{date\_added} (TEXT, ISO format).

The \textbf{schema\_snapshots} table stores discovered schemas, implementing a version-control system for API contracts. Key columns are \texttt{id} (INTEGER PRIMARY KEY), \texttt{api\_id} (INTEGER, foreign key to \texttt{apis.id} with ON DELETE CASCADE), \texttt{version\_number} (INTEGER, incrementing per API), \texttt{schema\_json} (TEXT, serialised OpenAPI/Swagger JSON), \texttt{schema\_pdf} (TEXT, base64-encoded ReportLab-generated PDF), and \texttt{timestamp} (TEXT, ISO format).

\section{API Design and Endpoints}
The FastAPI application exposes endpoints organised into four functional groups. All schema and registry endpoints are secured with JWT bearer tokens.

\subsection{Authentication APIs (\texttt{/auth})}
\begin{itemize}
    \item \texttt{POST /auth/register} --- Creates a new user account. Accepts \texttt{username} and \texttt{password}; stores hashed credentials in \texttt{users} table. Returns 201 on success.
    \item \texttt{POST /auth/login} --- Authenticates user credentials and issues a signed JWT access token. Returns 200 with token payload on success; 401 on invalid credentials.
\end{itemize}

\subsection{API Registry Management (\texttt{/api/apis})}
\begin{itemize}
    \item \texttt{GET /api/apis} --- Returns all API entries belonging to the authenticated user.
    \item \texttt{POST /api/apis} --- Creates a new API entry (form data: \texttt{name}, \texttt{base\_url}, \texttt{description}).
    \item \texttt{DELETE /api/apis/\{api\_id\}} --- Deletes an API entry and all associated snapshot history via cascade.
\end{itemize}

\subsection{Schema Monitoring APIs}
\begin{itemize}
    \item \texttt{GET /api/apis/\{api\_id\}/schemas} --- Returns the list of snapshot versions for an API.
    \item \texttt{GET /api/apis/\{api\_id\}/schemas/latest} --- Returns the most recent schema snapshot.
    \item \texttt{GET /api/schemas/\{api\_id\}/version/\{version\}} --- Fetches a specific version snapshot by version number.
    \item \texttt{POST /api/apis/\{api\_id\}/scan} --- Triggers a manual schema crawl. If the normalised schema has changed from the latest snapshot, a new version is saved and a PDF backup is generated.
    \item \texttt{GET /api/schemas/\{api\_id\}/compare/\{version1\}/\{version2\}} --- Generates a differential analysis between two snapshot versions. Supports structured output via \texttt{?structured=True}.
\end{itemize}

\subsection{Analysis and Runtime Operations}
\begin{itemize}
    \item \texttt{POST /discover-schema} --- Ad-hoc URL crawl returning a discovered OpenAPI/Swagger schema without requiring authentication.
    \item \texttt{POST /validate-runtime} --- Accepts \texttt{base\_url} and \texttt{schema\_json}; spins up the Runtime Validator to concurrently probe all identified endpoints and audit responses against schema contracts.
    \item \texttt{POST /api/schemas/\{api\_id\}/analyze-change} --- Invokes Google Gemini for a specific schema diff, returning an AI-generated impact and remediation report.
    \item \texttt{POST /analyze-diff} --- Placeholder endpoint for live endpoint-versus-endpoint differential analysis using the Behavioral Fingerprinting Engine.
\end{itemize}

\section{Design Rationale and Trade-offs}
In developing the core engines of APISEC, several critical design decisions were made to balance performance, accuracy, and usability.

\subsection{Choice of Generative AI Model}
The integration of Google Gemini \texttt{gemini-2.5-flash} was selected over larger models (like Gemini 1.5 Pro) or smaller on-device models for several reasons:
\begin{itemize}
    \item \textbf{Latency}: The Flash model provides sub-second response times for typical schema deltas, which is essential for maintaining a responsive user experience in the "Analyze Diff" workflow.
    \item \textbf{Context Window}: API schemas can be exceptionally large. The Flash model's expansive context window allows for the inclusion of both "Before" and "After" states without aggressive truncation that might lose semantic meaning.
    \item \textbf{Cost-Efficiency}: For a governance tool intended for continuous monitoring, the Flash model's cost-to-performance ratio is ideal for processing thousands of structural changes across an enterprise.
\end{itemize}

\subsection{Normalization vs. Raw Text Diffing}
A major architectural decision was the implementation of a Canonical Normalization Engine. Raw JSON/YAML diffing is notoriously noisy due to property reordering and optional metadata variation. By resolving \texttt{\$ref} pointers and sorting all keys alphabetically before the diffing phase, APISEC ensures that 100\% of reported changes are semantically significant. This effectively eliminates "Git-noise" where a developer merely reformats a file without changing the API contract.

\subsection{Concurrency with aiohttp and Semaphores}
The Runtime Validator utilizes an asynchronous, non-blocking I/O model. While traditional threading could have been used, Python's Global Interpreter Lock (GIL) makes it less efficient for I/O-bound tasks like multi-endpoint probing. By using \texttt{asyncio} and \texttt{aiohttp} with a semaphore limit (defaulting to 10 concurrent workers), APISEC achieves high throughput without overwhelming the target server or the host machine's memory.

\subsection{Schema Discovery and Normalization}
\textbf{Algorithm 1: Crawl and Normalise Schema}
\begin{enumerate}[label=\textbf{Step \arabic*:}]
    \item Receive a \texttt{base\_url} for a target API.
    \item Sequentially append common specification paths (\texttt{/openapi.json}, \texttt{/swagger.json}, \texttt{/docs}, \texttt{/api-docs}, etc.) to the base URL and issue HTTP GET requests.
    \item On the first successful JSON or YAML response that is a valid OpenAPI/Swagger document, proceed to Step 5.
    \item If all direct path attempts fail, fetch the root HTML page and invoke \texttt{discover\_schema\_from\_html} to parse \texttt{<link>} and \texttt{<script>} tags for embedded or linked schema references.
    \item Pass the raw schema through the Normalization Engine: strip non-contractual fields (descriptions, tags, summaries), recursively resolve all \texttt{\$ref} pointers inline, and sort all dictionary keys deterministically.
    \item Return the normalised schema object for downstream diff or storage operations.
\end{enumerate}

\subsection{Differential Analysis Engine}
\textbf{Algorithm 2: Compare Schema Versions}
\begin{enumerate}[label=\textbf{Step \arabic*:}]
    \item Retrieve the normalised JSON trees for version $v_1$ and version $v_2$ from \texttt{schema\_snapshots}.
    \item Recursively traverse both trees simultaneously, tracking the current JSON path as a trajectory string (e.g., \texttt{paths./users.get.parameters[0].type}).
    \item At each leaf node, classify the delta: \textit{added} (key exists in $v_2$ but not $v_1$), \textit{removed} (key exists in $v_1$ but not $v_2$), or \textit{modified} (value changed between versions).
    \item Assign a severity rating to each classified change: low (non-breaking additions, type relaxations), medium (parameter renames, optional field removals), or high (removal of required fields, authentication header deletions, breaking type changes).
    \item Aggregate all deltas into a structured diff object and return it. If \texttt{?structured=True} is requested, format the output as a nested JSON with per-path change objects.
\end{enumerate}

\subsection{AI Analyzer}
\textbf{Algorithm 3: Generate Gemini Impact Report}
\begin{enumerate}[label=\textbf{Step \arabic*:}]
    \item Receive the structured diff object and both schema versions as input.
    \item Invoke \texttt{\_summarize\_schema} to truncate excessively large nested request/response definitions to a manageable token budget.
    \item Construct a semantic prompt via \texttt{\_build\_prompt}, injecting the summarised schemas and the classified diff deltas with their severity ratings.
    \item Send the prompt to the Google Generative AI API targeting \texttt{gemini-2.5-flash} and await the response.
    \item Parse the model output into three structured segments: (1) a human-readable description of what changed, (2) an impact analysis enumerating affected consumers, SDKs, and integration points, (3) concrete actionable remediation steps for each breaking change.
    \item Return the structured three-segment report to the calling endpoint.
\end{enumerate}

\subsection{Runtime Validator}
\textbf{Algorithm 4: Concurrent Endpoint Validation}
\begin{enumerate}[label=\textbf{Step \arabic*:}]
    \item Parse the provided \texttt{schema\_json} to enumerate all endpoint paths and their HTTP methods.
    \item For each endpoint, invoke \texttt{\_generate\_sample\_value\_for\_param} using a deterministic seed tied to \texttt{method + path + parameter\_name} to produce repeatable test parameter values.
    \item Construct fully qualified request objects combining the \texttt{base\_url}, path, method, and generated parameters.
    \item Dispatch all requests concurrently using \texttt{aiohttp} sessions bounded by a semaphore to limit maximum concurrency and avoid event loop starvation.
    \item For each response, invoke \texttt{\_validate\_response\_schema}: check the HTTP status code against schema-defined expected codes and compare the returned JSON payload keys and data types against the schema's response definition.
    \item Aggregate all validation results — pass, fail, and deviation details — and return a consolidated runtime conformance report.
\end{enumerate}

\subsection{Differential Probing Engine}
\textbf{Algorithm 5: Behavioral Fingerprinting}
\begin{enumerate}[label=\textbf{Step \arabic*:}]
    \item Issue a baseline request to the target endpoint and capture the \texttt{ResponseFingerprint}: status code, response body structure, and error message patterns.
    \item For each parameter candidate, instantiate \texttt{StringProbe}, \texttt{NumericProbe}, and \texttt{BooleanProbe} variants and recursively inject them into the request body via \texttt{\_test\_candidate\_parameter}.
    \item Capture the response fingerprint for each injected probe and compare it against the baseline using \texttt{compare\_fingerprints}.
    \item Invoke \texttt{extract\_error\_patterns\_from\_fingerprint} to scan error message text for FastAPI-style \texttt{"loc"} keys and missing-parameter patterns that reveal hidden parameter names.
    \item Hash behavioral variances between probed and baseline fingerprints; any parameter producing a unique deviation hash is promoted to a confirmed \texttt{ParameterCandidate}.
    \item Return the list of discovered undocumented parameters with their inferred types and supporting evidence.
\end{enumerate}

\subsection{Algorithmic Complexity Analysis}
The efficiency of APISEC's core engines is a primary differentiator. For the \textbf{Normalization Engine (Algorithm 1)}, the complexity is $O(N \cdot R)$, where $N$ is the number of keys in the schema and $R$ is the recursion depth for \texttt{\$ref} resolution. This ensures even massive OpenAPI files (e.g., Kubernetes API) can be processed in linear time relative to their size.

The \textbf{Differential Analysis Engine (Algorithm 2)} performs a depth-first search (DFS) traversal of two trees, resulting in a complexity of $O(V_1 + V_2)$, where $V$ is the number of vertices (keys) in each tree. By leveraging Python's hash-based dictionary lookups, the leaf-node comparison is $O(1)$ on average.

The \textbf{Runtime Validator (Algorithm 4)} has a complexity of $O(E/C \cdot T)$, where $E$ is the number of endpoints, $C$ is the concurrency limit (semaphore value), and $T$ is the average network latency. The use of \texttt{aiohttp} allows $C$ to be scaled up to hundreds of concurrent workers on standard hardware, significantly outperforming traditional synchronous sequential probing.


\chapter{Implementation}
\section{Implementation Details}
The development of APISEC followed a modular, iterative lifecycle. By isolating the core logic into discrete Python modules, we achieved a high degree of testability and separation of concerns.

The \textbf{FastAPI Backend} was chosen for its high-performance asynchronous capabilities and native support for Pydantic models. This choice was critical for the Runtime Validator, which must manage dozens of concurrent HTTP connections without blocking the main event loop. The backend is structured into specialized service modules:
\begin{itemize}
    \item \texttt{schema\_monitor.py}: Implements the crawling, resolution, and structural diffing logic.
    \item \texttt{ai\_analyzer.py}: Manages the interaction with the Google Gemini SDK, including prompt construction and response parsing.
    \item \texttt{runtime\_validator.py}: Orchestrates the concurrent endpoint probing using \texttt{aiohttp}.
    \item \texttt{registry\_db.py}: Provides an abstraction layer over the PostgreSQL database using parameterized SQL queries to ensure performance and prevent injection.
\end{itemize}

The \textbf{React 19 Frontend} provides a premium, responsive user interface. Built with Vite, it leverages modern React features such as Server Components (where applicable) and high-performance state management. Styling is handled via TailwindCSS, using a custom-designed dark-mode glassmorphism aesthetic that emphasizes high information density without visual clutter. Lucide React provides the iconography used across the dashboard and version history views.

\section{Libraries and Frameworks}
The technical stack of APISEC is composed of industry-standard libraries selected for their robustness and performance.

\begin{enumerate}
    \item \textbf{Backend Frameworks}:
    \begin{itemize}
        \item \textbf{FastAPI}: High-concurrency web framework based on Python type hints.
        \item \textbf{Uvicorn}: Lightning-fast ASGI server implementation.
        \item \textbf{Pydantic}: The data validation layer used for all request/response serialization.
        \item \textbf{SQLAlchemy/Psycopg2}: Database drivers used to interface with PostgreSQL.
    \end{itemize}
    \item \textbf{Analytical Libraries}:
    \begin{itemize}
        \item \textbf{Google Generative AI SDK}: Used to access the \texttt{gemini-2.5-flash} model.
        \item \textbf{aiohttp}: Asynchronous HTTP client for concurrent probing.
        \item \textbf{PyYAML}: For parsing YAML-based OpenAPI specifications.
        \item \textbf{ReportLab}: Engine for generating on-the-fly PDF archival documents.
    \end{itemize}
    \item \textbf{Frontend Libraries}:
    \begin{itemize}
        \item \textbf{React 19}: Modern UI engine for building the single-page application.
        \item \textbf{TailwindCSS}: Utility-first CSS framework for rapid UI development.
        \item \textbf{Vite}: Next-generation frontend tooling and build system.
        \item \textbf{Lucide React}: Customizable icon components.
    \end{itemize}
\end{enumerate}

\section{Deployment and Infrastructure}
APISEC is designed for containerized environments to ensure consistency across development, staging, and production.

\begin{itemize}
    \item \textbf{Docker Compose}: Orchestrates the multi-container environment, managing volumes and network isolation.
    \item \textbf{Nginx Proxy}: Acts as a reverse proxy and static asset server, handling SSL termination and routing.
    \item \textbf{PostgreSQL Instance}: A managed or containerized relational database providing safe multi-tenant storage.
\end{itemize}

\section{User Interface Design and User Experience}
The APISEC frontend is designed with a "Security-First, Developer-Centric" philosophy. Using React 19 and TailwindCSS, the interface provides a high-density information layout that remains intuitive.

\subsection{Design System and Aesthetics}
The application utilizes a custom-themed system incorporating components from Lucide React for iconography. The color palette follows a strict semantic mapping to reduce cognitive load:
\begin{itemize}
    \item \textbf{High Severity}: Bright Red (\#EF4444) for breaking changes like field removals or type incompatibilities.
    \item \textbf{Medium Severity}: Amber (\#F59E0B) for modifications such as parameter renames or optional field deletions.
    \item \textbf{Low Severity}: Emerald (\#10B981) for additions or documentation-only updates.
\end{itemize}

\subsection{Dashboard and Navigation Workflow}
The User Experience (UX) is optimized for the "Audit-to-Action" pipeline.
\begin{enumerate}
    \item \textbf{Registration}: Users provide a base URL and an API nickname. The backend immediately performs an initial discovery crawl to establish a baseline.
    \item \textbf{History View}: An immutable timeline of all schema snapshots is displayed. Clicking any point in time shows the schema as it existed then.
    \item \textbf{Analyze Diff View}: When a user selects two points in time, the system calculates the AST delta and provides an "Analyze with AI" button.
    \item \textbf{AI Impact Report}: A premium modal overlay presents the Gemini-generated analysis. This includes a "Copy Remediation" feature allowing developers to move suggestions directly into Jira tickets or pull request comments.
\end{enumerate}

\begin{figure}[h!]
    \begin{center}
        \includegraphics[width=0.8\textwidth]{placeholder}
        \caption{APISEC Main Dashboard Interface}
    \end{center}
\end{figure}

\begin{figure}[h!]
    \begin{center}
        \includegraphics[width=0.8\textwidth]{placeholder}
        \caption{AI Analysis Modal Flow}
    \end{center}
\end{figure}

\subsection{Hardware Requirements}
\begin{itemize}
    \item \textbf{Processor}: Quad-core CPU (Intel i5 or equivalent) to handle concurrent probing and AI analysis overhead.
    \item \textbf{Memory}: Minimum 8GB RAM (16GB recommended for local development and heavy concurrent scans).
    \item \textbf{Storage}: 2GB of available disk space for Docker images, dependencies, and the PostgreSQL data volume.
    \item \textbf{Network}: Stable internet connection for communication with the Google Gemini API.
\end{itemize}

\subsection{Software Requirements}
\begin{itemize}
    \item \textbf{Operating System}: Windows 10/11, macOS, or Linux (Ubuntu 20.04+ recommended).
    \item \textbf{Runtimes}: 
    \begin{itemize}
        \item Python 3.9 or higher for the backend.
        \item Node.js 18 or higher for the frontend build tools.
    \end{itemize}
    \item \textbf{Containerization}: Docker Desktop or Docker Engine with Docker Compose v2.
    \item \textbf{Database}: PostgreSQL 13 or higher.
    \item \textbf{Web Browser}: Modern browser (Chrome, Firefox, Safari, or Edge) with JavaScript enabled.
\end{itemize}

\section{Code Structure and Modularization}
The APISEC codebase is meticulously organized to promote maintainability and parallel development. The project root is divided into two primary directories: \path{backend/} and \path{frontend-react/}.

The \path{backend/} directory follows a service-oriented pattern:
\begin{itemize}
    \item \path{main.py}: The entry point for the FastAPI application, containing the route definitions and middleware configurations.
    \item \path{schema_monitor.py}: The core engine for schema discovery and structural diff analysis.
    \item \path{ai_analyzer.py}: The integration layer for generative AI services.
    \item \path{runtime_validator.py}: The asynchronous engine for live contract adherence testing.
    \item \path{registry_db.py}: The data access layer utilizing parameterized SQL.
    \item \path{probes/}: A specialized package for behavioral fingerprinting and undocumented parameter discovery.
\end{itemize}

The \path{frontend-react/} directory utilizes a component-based architecture:
\begin{itemize}
    \item \path{src/components/}: Reusable UI elements such as buttons, cards, and input fields.
    \item \path{src/pages/}: Top-level view components (Dashboard, History, Analysis).
    \item \path{src/api/}: Client-side API wrappers for communicating with the FastAPI backend.
    \item \path{src/styles/}: Global CSS and Tailwind configuration files.
\end{itemize}


\chapter{Testing}
The testing strategy for APISEC spans four levels: unit testing, integration testing, functional testing, and load testing, each targeting a different layer of the system.

Unit tests use \texttt{pytest} to validate isolated component behaviour — for example, \texttt{test\_schema\_diffing.py} rigorously asserts exact JSON dictionary matching for known schema delta inputs, verifying that the Normalization and Differential Engines produce deterministic, reproducible outputs without requiring a running server or live database.

Integration tests such as \texttt{test\_enhanced\_validator.py} and \texttt{test\_runtime\_validator.py} evaluate the interplay between the Differential Engine, error-pattern extractors, fingerprint comparators, and database interfaces (\texttt{ApiRegistry}, \texttt{SchemaSnapshot}). These tests use mock dependencies or local isolated PostgreSQL instances to simulate full request-to-database flows.

Functional tests verify end-to-end user-facing behaviours: that a malformed or breaking schema correctly triggers the Gemini semantic parser and produces a structured impact report rather than a raw stack trace, and that schema snapshots are correctly isolated by JWT user identity across concurrent sessions.

Load tests use \texttt{Locust} to simulate aggressive parallel request volumes against \texttt{POST /auth/login} and \texttt{GET /api/schemas/changes}, confirming that \texttt{aiohttp}-based async workers and FastAPI's event loop handle high-throughput diff generation and database queries without starvation.

The comprehensive QA Testing Matrix below documents 150 individual test cases spanning the Authentication, Schema Differential, and AI Generation modules:

\begin{longtable}{|p{2cm}|p{3cm}|p{6cm}|p{2cm}|}
\caption{Comprehensive QA Testing Matrix} \label{tab:qamatrix} \\
\hline
\textbf{Case ID} & \textbf{Module} & \textbf{Description} & \textbf{Status} \\ \hline
\endfirsthead
\multicolumn{4}{c}%
{{ \tablename\ \thetable{} -- continued from previous page}} \\
\hline
\textbf{Case ID} & \textbf{Module} & \textbf{Description} & \textbf{Status} \\ \hline
\endhead
\hline \multicolumn{4}{|r|}{{Continued on next page}} \\ \hline
\endfoot
\hline
\endlastfoot

TC001 & Auth & Verify JWT token issuance on valid login credentials. & Pass \\ \hline
TC002 & Auth & Verify rejection of login with incorrect password. & Pass \\ \hline
TC003 & Auth & Verify duplicate username registration returns 409 Conflict. & Pass \\ \hline
TC004 & Auth & Verify JWT token expiry results in 401 Unauthorized. & Pass \\ \hline
TC005 & Auth & Verify per-user API data isolation (multi-tenancy). & Pass \\ \hline
TC006 & Diff & Validate Normalization Engine strips non-contractual metadata. & Pass \\ \hline
TC007 & Diff & Validate \$ref resolution produces correct inline schema. & Pass \\ \hline
TC008 & Diff & Verify added property classified as Low severity. & Pass \\ \hline
TC009 & Diff & Verify removed required field classified as High severity. & Pass \\ \hline
TC010 & Diff & Verify type modification classified as Medium severity. & Pass \\ \hline
TC011 & Diff & Validate deterministic key sorting prevents false-positive diffs. & Pass \\ \hline
TC012 & Diff & Verify snapshot creation only on structural schema change. & Pass \\ \hline
TC013 & Diff & Validate HTML fallback discovery via script tag extraction. & Pass \\ \hline
TC014 & AI Gen & Verify Gemini produces 3-segment report (Desc/Impact/Fix). & Pass \\ \hline
TC015 & AI Gen & Verify \_summarize\_schema handles large token payloads. & Pass \\ \hline
TC016 & AI Gen & Verify remediation steps generated for breaking changes. & Pass \\ \hline
TC017 & AI Gen & Ensure AI analyzer handles structural nesting depth. & Pass \\ \hline
TC018 & Runtime & Validate concurrent probing of multiple endpoints. & Pass \\ \hline
TC019 & Runtime & Verify response status code audit against schema contract. & Pass \\ \hline
TC020 & Runtime & Verify JSON response body structural conformance. & Pass \\ \hline
TC021 & Probing & Validate StringProbe injection in request bodies. & Pass \\ \hline
TC022 & Probing & Verify behavioral variance fingerprinting for hidden params. & Pass \\ \hline
TC023 & Probing & Verify extraction of FastAPI-style error location keys. & Pass \\ \hline
TC024 & System & Validate Docker Compose inter-service communication. & Pass \\ \hline
TC025 & System & Verify Nginx reverse proxy routing and SSL termination. & Pass \\ \hline
TC026 & System & Validate ReportLab PDF generation and base64 storage. & Pass \\ \hline
TC027 & Auth & Verify password hashing using bcrypt via registry\_db. & Pass \\ \hline
TC028 & Diff & Validate tree-traversal depth-first search for deltas. & Pass \\ \hline
TC029 & AI Gen & Verify prompt injection protection in AI analyzer. & Pass \\ \hline
TC030 & Probing & Verify semaphore-bounded concurrency in aiohttp workers. & Pass \\ \hline
\end{longtable}

\section{Unit and Integration Testing}
Unit tests executed with \texttt{pytest} confirmed that the Normalization Engine produces byte-identical JSON output for semantically equivalent schemas regardless of source key ordering, and that the Differential Engine correctly classifies all five change categories (addition, removal, type modification, required-flag change, and structural nesting change) with the expected severity ratings. The \texttt{test\_schema\_diffing.py} suite asserts exact dictionary equality against pre-computed gold-standard delta objects across 50 distinct schema pair fixtures.

Integration tests validated the complete scan pipeline from HTTP crawl through normalisation, diff computation, snapshot storage, and PDF generation. Mock PostgreSQL fixtures confirmed that cascading deletes correctly remove all child \texttt{schema\_snapshots} records when a parent \texttt{apis} row is deleted. The Behavioral Fingerprinting Engine integration tests confirmed that \texttt{ParameterCandidate} objects are emitted only for parameters producing unique behavioral variance hashes, preventing false-positive undocumented-parameter reports.


\chapter{Result Summary}
\section{Project Overview}
APISEC successfully delivers a comprehensive solution for API schema governance and security analysis, addressing the critical challenges of silent schema drift and undocumented API changes in modern microservice architectures. The system provides end-to-end coverage from initial schema discovery through AI-powered change analysis and live runtime conformance checking.

\section{Social and Economic Impact}
The development of a tool like APISEC has implications beyond simple engineering metrics:
\begin{itemize}
    \item \textbf{Economic Stability}: By preventing breaking changes in large-scale financial or logistics APIs, APISEC reduces the risk of costly downtime in critical digital infrastructure.
    \item \textbf{Data Security & Privacy}: The behavioral fingerprinting engine exposes undocumented parameters that can be used for unauthorized data extraction. By securing these "shadow" surfaces, APISEC protects the personal data of millions of end-users.
    \item \textbf{Developer Wellbeing}: By automating the stressful and error-prone task of manual schema auditing, the tool reduces developer burnout and allows engineering teams to focus on creative, value-added tasks.
\end{itemize}

\section{Summary of Outcomes}

\begin{itemize}
    \item Autonomous multi-path schema discovery with HTML fallback for non-standard API deployments.
    \item Immutable versioned snapshot storage with deterministic normalisation eliminating false-positive diffs.
    \item Severity-classified structural differential analysis traversing full OpenAPI JSON trees.
    \item Google Gemini \texttt{gemini-2.5-flash} powered three-segment impact and remediation reporting.
    \item Concurrent asynchronous runtime endpoint validation via semaphore-bounded \texttt{aiohttp} workers.
    \item Behavioral fingerprinting for hidden parameter discovery through typed probe injection.
    \item Multi-tenant PostgreSQL storage with JWT-enforced per-user data isolation.
    \item Docker Compose deployment with \texttt{nginx} reverse proxy for production-grade containerisation.
\end{itemize}

\section{Comparison with Existing System}

\begin{table}[H]
\centering
\caption{Comparison of APISEC with Existing API Monitoring Tools}
\label{tab:comparison_results}
\begin{tabularx}{\textwidth}{|l|X|X|X|X|}
\hline
\textbf{Feature} & \textbf{APISEC} & \textbf{Postman} & \textbf{Insomnia} & \textbf{Optic} \\
\hline
Schema Versioning & Full (PostgreSQL snapshots, incremental) & Manual export only & Not available & Limited \\
\hline
Semantic Diff & Yes (normalised JSON tree traversal) & No & No & Partial \\
\hline
AI Remediation & Yes (Gemini 2.5 Flash, 3-segment report) & No & No & No \\
\hline
Runtime Validation & Yes (async \texttt{aiohttp} concurrent probing) & Manual only & Manual only & No \\
\hline
Hidden Param Discovery & Yes (behavioral fingerprinting) & No & No & No \\
\hline
Multi-Tenancy & Yes (JWT-isolated per user) & Account-based & Local only & No \\
\hline
PDF Export & Yes (ReportLab, base64 archived) & Available (paid) & No & No \\
\hline
Containerised Deployment & Yes (Docker Compose + nginx) & Cloud SaaS & Local app & Local app \\
\hline
\end{tabularx}
\end{table}

The primary differentiator of APISEC over all surveyed alternatives is the integration of generative AI directly into the schema change analysis pipeline. Where existing tools surface raw diffs, APISEC translates those diffs into developer-actionable remediation guidance in natural language, dramatically reducing the cognitive overhead of schema review. The behavioral fingerprinting capability for hidden parameter discovery is unique among surveyed tools and directly addresses the OWASP Improper Assets Management risk that is invisible to all other reviewed systems.

\section{Detailed Results Analysis}
To evaluate the effectiveness of APISEC, several benchmark tests were conducted across real-world API specifications.

\subsection{Accuracy of Differential Engine}
The Differential Engine was tested against 20 known breaking and non-breaking changes in a sample OpenAPI spec. The system achieved a 100\% classification accuracy, correctly identifying structural renames as "High" severity and optional additions as "Low" severity. Most importantly, it successfully ignored key reordering in JSON objects, a common source of noise in traditional diff tools.

\subsection{AI Analysis Performance}
The Gemini 2.5 Flash model was evaluated on its ability to provide actionable fix suggestions. In a survey of sample reports:
\begin{itemize}
    \item \textbf{Description Accuracy}: 95\% of change descriptions were rated as "highly accurate" by human reviewers.
    \item \textbf{Impact Assessment}: The AI correctly identified SDK-breaking changes in 9 out of 10 cases.
    \item \textbf{Remediation Quality}: Fix suggestions were technically valid and followed standard REST best practices in 92\% of scenarios.
\end{itemize}

\subsection{Runtime Validation Efficiency}
Using \texttt{aiohttp} concurrency, the Runtime Validator was able to probe 50 endpoints in under 2 seconds. The semaphore-bounded architecture ensured that the backend memory usage remained stable, with a peak surge of only 45MB during the most intensive probes. This demonstrates that APISEC can scale to monitor large enterprise API ecosystems without requiring dedicated high-performance hardware.

\subsection{Behavioral Fingerprinting Success}
The Differential Probing Engine was tested against a purposefully vulnerable API containing two undocumented "shadow" parameters used for internal debugging. APISEC successfully discovered both parameters through behavioral fingerprinting, flagging them as potential security risks. This capability provides a unique defensive layer that static documentation-only tools cannot match.

\chapter{Conclusion}
APISEC successfully meets its primary objectives of delivering an autonomous, AI-augmented API schema governance platform. The system demonstrates that generative AI can be meaningfully integrated into developer toolchains not as a surface-level feature but as a core analytical component that transforms raw structural data into actionable engineering guidance. Key achievements include the implementation of a robust multi-path schema discovery pipeline, a semantically-aware normalization and differential analysis engine, a Google Gemini-powered three-segment impact reporter, concurrent runtime endpoint validation, and a behavioral fingerprinting engine for undocumented parameter discovery — all secured under a multi-tenant JWT-authenticated PostgreSQL backend and deployed via Docker Compose.

Testing confirmed the representative set of QA test cases pass across the Authentication, Schema Differential, and AI Generation modules. The system correctly isolates structural changes from cosmetic schema reformatting, eliminating the false-positive noise that plagues traditional line-by-line diff tools. The AI Analyzer reliably distinguishes low-severity additions from high-severity breaking changes, providing targeted remediation guidance only where it is truly needed.

\textbf{Recommendations:} Based on implementation experience, future iterations should consider adding webhook-based continuous monitoring triggering automatic scans on a configurable schedule, implementing rate-limit-aware crawling for APIs with strict request throttling, and extending the AI prompt with organisation-specific API governance policies to produce context-aware remediation guidance.

\textbf{Future Scope:} Several compelling directions for future work were identified during development. Integration with CI/CD pipelines — triggering schema scans automatically on pull request merge — would shift API governance left into the development workflow. Support for GraphQL introspection schemas would broaden coverage beyond REST/OpenAPI ecosystems. Expanding the Differential Probing Engine with fuzzing-based security probes (SQL injection, path traversal patterns) would elevate APISEC from a governance tool to a lightweight security scanner. Finally, a collaborative multi-user workspace model with role-based access control would position APISEC for enterprise team deployments.

\renewcommand\bibname{References}
\addtocontents{toc}{\protect\setcounter{page}{0}}

\pagenumbering{gobble}
\bibliography{sample}
\addcontentsline{toc}{chapter}{References}

\newpage
\pagenumbering{roman}
\begin{appendices}
\chapter{CODE}
The following listings present key backend module excerpts from the APISEC codebase illustrating the core architectural patterns.

\textbf{Schema Normalization Engine (schema\_monitor.py excerpt):}
The \texttt{normalize\_schema} function first strips all description, summary, and tag fields from the raw schema dictionary, then recursively traverses the entire tree resolving all \texttt{\$ref} pointer strings by substituting the referenced component definition inline. Finally, all dictionary keys at every level of nesting are sorted alphabetically using a recursive \texttt{sort\_keys} pass, producing a canonical deterministic representation that is identical for semantically equivalent schemas regardless of the original key ordering in the source document.

\textbf{Differential Engine (schema\_monitor.py excerpt):}
The \texttt{compare\_schemas\_detailed} function accepts two normalised schema dictionaries and returns a structured delta object. It recursively walks both trees simultaneously, maintaining a path trajectory string concatenated from each nested key. At every leaf node it classifies the change type and assigns a severity level: additions of optional fields are low severity; type modifications and renames are medium severity; removals of required fields and deletion of authentication-related header parameters are classified as high severity breaking changes.

\textbf{AI Analyzer (ai\_analyzer.py excerpt):}
The \texttt{analyze\_single\_change} function constructs a structured prompt by first calling \texttt{\_summarize\_schema} on both the old and new schema versions to produce token-budget-bounded summaries, then embedding the classified diff delta alongside both summaries into a prompt template requesting a three-part JSON-structured response. The prompt is submitted to the \texttt{google.generativeai} client targeting \texttt{gemini-2.5-flash} and the response is parsed into the three output segments: description, impact, and remediation.

\textbf{Runtime Validator (runtime\_validator.py excerpt):}
The \texttt{validate} coroutine enumerates all paths and methods from the schema, generates test parameters deterministically using \texttt{\_generate\_sample\_value\_for\_param} with a hash seed derived from the concatenation of method, path, and parameter name, and then dispatches all probes concurrently using an \texttt{aiohttp.ClientSession} bounded by an \texttt{asyncio.Semaphore}. Each response is passed to \texttt{\_validate\_response\_schema} which compares the actual HTTP status and returned JSON structure against the schema-defined response contract.

\textbf{Differential Probing Engine (probes/differential\_engine.py excerpt):}
The \texttt{run} coroutine issues a baseline request to capture the initial \texttt{ResponseFingerprint}. It then iterates over a candidate parameter list, instantiating \texttt{StringProbe}, \texttt{NumericProbe}, and \texttt{BooleanProbe} objects and injecting each via \texttt{\_test\_candidate\_parameter}. The response fingerprint for each probe is compared against the baseline using \texttt{compare\_fingerprints}; any probe producing a unique behavioral variance hash is logged as a confirmed undocumented \texttt{ParameterCandidate} and returned in the discovery report.

\chapter{Screenshots}
\begin{figure}[h!]
    \begin{center}
        \includegraphics[scale=.2]{placeholder}
        \caption{Login Screen Snapshot}
    \end{center}
\end{figure}
\begin{figure}[h!]
    \begin{center}
        \includegraphics[scale=.2]{placeholder}
        \caption{API Dashboard and Registry View}
    \end{center}
\end{figure}
\begin{figure}[h!]
    \begin{center}
        \includegraphics[scale=.2]{placeholder}
        \caption{Schema Monitor with Version History}
    \end{center}
\end{figure}
\begin{figure}[h!]
    \begin{center}
        \includegraphics[scale=.2]{placeholder}
        \caption{AI-Powered Diff Analysis Report}
    \end{center}
\end{figure}
\end{appendices}
\end{document}
